\section{Marco Teórico}
\label{sec:marco-teorico}

El marco teórico revisa cuatro trabajos que sostienen, en conjunto, las decisiones técnicas
del prototipo: dos resuelven el problema de cuántos y dónde (conteo y localización), y dos
resuelven el problema de quién sigue siendo quién a lo largo del tiempo (seguimiento).

\subsection{CSRNet: redes convolucionales dilatadas para conteo de multitudes (Li, Zhang
y Chen, 2018)}
\label{subsec:csrnet}

\textbf{Qué propone.} Una red completamente convolucional con un \emph{front-end} que
reutiliza las primeras diez capas de VGG-16 y un \emph{back-end} de convoluciones dilatadas
que amplía el campo receptivo sin \emph{pooling}. La salida es un mapa de densidad continuo,
no un número ni una posición individual.

\textbf{Resultados.} Sobre ShanghaiTech, reduce el MAE en 7\,\% respecto de CP-CNN en la
Parte A y en 47.3\,\% en la Parte B; sobre TRANCOS, 15.4\,\% menos que el mejor método
previo.

\textbf{Por qué aplica al caso LAP.} Su salida es agregada por construcción, no requiere
rasgos faciales; es puramente convolucional y admite resoluciones de entrada variables,
relevante porque las cámaras de un terminal no comparten resolución, altura ni ángulo.

\textbf{Limitación reconocida por los autores.} En escenas dispersas (UCSD), CSRNet
queda por detrás de la mayoría de métodos salvo MCNN en MAE: su ventaja se concentra en
alta densidad, y buena parte de la jornada de un terminal aéreo transcurre en baja densidad.

\subsection{P2PNet: localización pura por puntos (Song et al., 2021)}
\label{subsec:p2pnet}

\textbf{Qué propone.} Un marco que predice directamente un conjunto de puntos con
coordenadas y confianza, sin representación intermedia, con emparejamiento uno a uno vía
algoritmo húngaro; aporta además la métrica \emph{density Normalized Average Precision}
(nAP).

\textbf{Resultados.} Sobre ShanghaiTech Parte A, reduce el MAE 4.8\,\% y el MSE 12.9\,\%
frente a ADSCNet; sobre UCF\_CC\_50 rebaja el MAE en 2.1 puntos; MAE de 85.32 sobre
UCF-QNRF, con nAP a 0.5 cercano al 90\,\% en la mayoría de los conjuntos.

\textbf{Por qué aplica al caso LAP.} La localización individual convierte una cifra en una
decisión operativa, pero eleva el perfil de riesgo normativo: un punto persistido y encadenado
en el tiempo se acerca a un seguimiento de trayectorias, una finalidad distinta de estimar
aglomeraciones.

\textbf{Limitación reconocida por los autores.} En UCF-QNRF su precisión no es tan
competitiva frente a ADSCNet; en NWPU-Crowd su desempeño es inferior por depender de un
único mapa de características de una sola escala.

\subsection{P2R: pérdida punto a región para conteo semi-supervisado (Lin, Zhao y Chan,
2025)}
\label{subsec:p2r}

\textbf{Qué propone.} Un marco semi-supervisado maestro-estudiante que sustituye el
emparejamiento punto a punto por uno punto a región, eliminando la necesidad del algoritmo
húngaro en esa etapa.

\textbf{Resultados.} Con solo 5\,\% de datos anotados, supera al aprendizaje totalmente
supervisado con 10\,\% (MAE 93.7, MSE 155.2 sobre ShanghaiTech Parte A); bajo supervisión
completa, supera a P2PNet en los cuatro conjuntos evaluados, con un cómputo de pérdida
$\approx$68 veces más rápido por imagen.

\textbf{Por qué aplica al caso LAP.} Conecta con el principio de minimización de datos de la
Ley 29733: un método competitivo con 5\,\% de datos etiquetados reduce en proporción similar
el volumen de imágenes que deberían anotarse manualmente.

\textbf{Limitación reconocida por los autores.} Fallo por desenfoque de movimiento de la
multitud, la norma y no la excepción en un terminal aéreo; los autores documentan
explícitamente la brecha de dominio.

\subsection{ByteTrack: seguimiento multiobjeto asociando toda caja de detección (Zhang
et al., 2022)}
\label{subsec:bytetrack}

\textbf{Qué propone.} Asociar todas las cajas de detección, no solo las de alta confianza:
para las de baja puntuación, aprovecha su similitud con las trayectorias activas para recuperar
objetos reales y descartar detecciones de fondo, combinado con un filtro de Kalman.

\textbf{Resultados.} Desempeño de referencia en los benchmarks MOT20, HiEve y
BDD100K.

\textbf{Por qué aplica al caso LAP.} El prototipo adapta esta lógica a puntos en el
seguimiento mono-cámara (Kalman de velocidad constante más asignación húngara en dos
etapas), y la extensión de monitoreo integrado usa ByteTrack sin modificar sobre cajas de un
detector YOLO11n.

\subsection{Síntesis de antecedentes}
\label{subsec:sintesis-antecedentes}

CSRNet resuelve cuántos con una superficie de densidad de la que no puede extraerse
dónde. P2PNet resuelve dónde eliminando la representación intermedia, pero paga ese avance
con un costo de anotación creciente. P2R resuelve ese costo llevando el enfoque a régimen
semi-supervisado. ByteTrack resuelve quién sigue siendo quién entre cuadros consecutivos sin
descartar evidencia débil. De la lectura conjunta se derivan tres criterios que gobiernan el
diseño del sistema: la granularidad de la salida es una decisión de cumplimiento normativo y
no solo de producto; la eficiencia de etiquetado de P2R es, además de una ventaja técnica, un
argumento de minimización de datos; y ninguna cifra de la literatura revisada es transferible al
Aeropuerto Jorge Chávez sin una validación local que hoy no se ha realizado.
